\chapter{Introducción}
\pagenumbering{arabic}
\label{chap:introtesis}

La economía de Honduras se desenvuelve en una compleja red de
interdependencias, donde el crecimiento del \textbf{Producto Interno
Bruto (PIB)} se entrelaza con la dinámica de la \textbf{inflación},
la \textbf{tasa de desempleo} y la \textbf{balanza comercial}
\cite{EconHonduras, WorldBankGDP}.  Estos indicadores no operan en
el vacío; su comportamiento está condicionado por la sostenibilidad
de la \textbf{deuda pública}, la capacidad de atraer \textbf{Inversión
Extranjera Directa (IED)} y la estabilidad de la \textbf{tasa de
cambio} \cite{BCH2023, WorldBank2023}.  En última instancia, el
resultado de esta interacción se refleja en el \textbf{Índice de
Desarrollo Humano (IDH)}, que encapsula la calidad de vida de la
población.  La volatilidad inherente a una economía emergente,
sumada a la incertidumbre política y a shocks externos, convierte
el análisis riguroso de estas variables en una tarea fundamental
para la planificación estratégica.

Navegar esta complejidad requiere más que intuición; exige
herramientas matemáticas capaces de modelar la incertidumbre y
proyectar escenarios futuros con solidez.  El análisis de series
temporales se vuelve aquí indispensable.  Sin embargo, la
construcción de pronósticos fiables en el contexto hondureño
tropieza con desafíos significativos: desde la dispersión de los
datos hasta la desconexión entre el análisis cuantitativo y la
toma de decisiones, lo que a menudo deja un vacío que se llena con
previsiones puramente cualitativas.

\section{Series de Estudio}

Para anclar el análisis en la realidad económica del país, esta
investigación se centra en dos familias de series temporales con
características deliberadamente distintas
(Capítulo~\ref{chap:seleccionvariable}).  Por un lado, el
\textbf{PIB per cápita en lempiras a precios constantes}, indicador
macroeconómico canónico de baja frecuencia cuyas anomalías
estructurales reflejan crisis históricas (2009, 2020).  Por otro,
los \textbf{precios semanales de combustibles} (Súper, Regular,
Diésel y Kerosene), series de alta frecuencia que capturan la
volatilidad del mercado energético para el período
\textbf{2017--2025}, recolectadas mediante \textit{web scraping}
automatizado del portal Proceso Digital \cite{ProcesoHN2025}.

\section{Marco Metodológico}

Frente a este panorama, la tesis desarrolla un marco metodológico
progresivo organizado en seis bloques temáticos:

\begin{enumerate}
    \item \textbf{Operador TCROC y variantes}
    (Capítulo~\ref{cap:tcroc}).  Se formaliza la familia de
    transformaciones TCROC, TCROCM, ETCROC y ETCROCM, demostrando
    existencia, unicidad, consistencia asintótica y estabilidad
    ante perturbaciones del estimador generalizado.

    \item \textbf{Integración estocástica}
    (Capítulo~\ref{cap:TCROC-Markov}).  Se conecta el operador TCROC
    con cadenas de Markov, estableciendo los fundamentos de teoría
    ergódica, $\phi$-mixing y estimación asintótica que sustentan la
    predicción estocástica de regímenes.

    \item \textbf{Modelo híbrido NNLS}
    (Capítulo~\ref{cap:tcroc_hibrido_integrado}).  Se presenta el estimador de
    Mínimos Cuadrados No Negativos para la matriz de transición,
    garantizando por construcción la validez estocástica de los
    parámetros.

    \item \textbf{Validación empírica en combustibles}
    (Capítulo~\ref{cap:regimenes_combustibles}).  Se aplica el marco completo a
    los precios semanales de combustibles hondureños, incluyendo
    optimización de hiperparámetros, análisis espectral y comparación
    con benchmarks (cuantiles fijos, MS-AR).

    \item \textbf{Extensión no lineal}
    (Capítulo~\ref{cap:sec_ssrc}).  Se propone la Computación de
    Reservorio Estructurado (SSRC) como generalización no lineal
    del marco lineal TCROC-Markov, evaluando su capacidad predictiva
    sobre las mismas series.

    \item \textbf{Conclusiones y trabajo futuro}
    (Capítulo~\ref{cap:conclusiones}).  Se sintetizan los hallazgos,
    las contribuciones y las líneas de investigación abiertas.
\end{enumerate}

Los pronósticos resultantes, aunque no son certezas, proveen un mapa
de escenarios posibles fundamental para la planificación estratégica
en contextos complejos y volátiles.
